% Modernised reading — fonds Alexandre Grothendieck, shelfmark 151, batch 4
% (pages 61-74).
%
% This is an INTERPRETATION, not a transcription. It re-reads the pages in
% today's notation and vocabulary, states the hypotheses the manuscript leaves
% implicit, and drops the critical apparatus so that the mathematics reads
% straight through. Everything the brackets and underlines carried — a doubtful
% reading, a notation he used differently, a missing passage — moves into a
% footnote.
%
% It is held to being mathematically correct as it stands, not merely faithful
% to the page. Where the manuscript is loose or mistaken, the true statement is
% what appears here, and the footnote says what the page has. In any dispute of
% substance the transcription governs: whoever cites Grothendieck must cite
% batch-04.fr.tex.
%
% Scope: the folder was read WHOLE before any of this was written — all four
% batches, pages 1-74, in one pass — and one file is written per batch, this
% being the last. The folder BREAKS OFF mid-construction on page 74, in the
% middle of section 4's part C, and this reading says so and supplies no
% ending. Four of its pages (62, 66, 68, 70) are transcribed as skeletons
% because the manuscript is written over to the point of illegibility; where
% this reading fills such a gap it says in a footnote that it is doing so.
%
% Pass: Opus 5 (claude-opus-5), 2026-08-16 — first pass, unchecked against the
% pages by a human.
\documentclass[11pt,a4paper]{article}
% Le préambule commun à toutes les transcriptions du fonds Grothendieck.
%
% Il tient en une idée : **l'incertitude se compose, elle ne se résout pas.**
% Une transcription qui lisse un mot douteux a détruit la seule chose pour
% laquelle on la faisait — un lecteur doit pouvoir savoir, à chaque endroit,
% ce qui était sur le feuillet et ce qui a été ajouté. Les six macros qui
% suivent sont donc l'appareil critique, et non de la décoration.
%
% Les mêmes commandes sont reconnues par `scripts/render.mjs`, qui produit la
% vue de lecture du site. Une macro ajoutée ici doit l'être aussi là-bas,
% faute de quoi le rendu échouera bruyamment — ce qui est voulu : un rendu
% silencieusement faux est pire qu'un rendu absent.

\ProvidesPackage{grothendieck}[2026/08/08 Transcriptions du fonds Grothendieck]

\usepackage{amsmath,amssymb,amsthm}
\usepackage{mathrsfs}
\usepackage{stmaryrd}
% Les diagrammes commutatifs sont partout dans ce fonds : c'est la langue
% même de la géométrie algébrique de Grothendieck, et une transcription qui
% les rendrait en prose ne transcrirait rien.
\usepackage{tikz-cd}
% Français par défaut — c'est la langue des feuillets — mais l'anglais est
% chargé aussi : la traduction et le résumé sont des documents anglais, et la
% typographie française (espaces avant les deux-points) y serait une faute.
% Ces documents basculent par \selectlanguage{english} en tête de corps.
\usepackage[english,french]{babel}
\usepackage{fontspec}
\usepackage{geometry}
\geometry{a4paper,margin=25mm,marginparwidth=28mm}
\usepackage{marginnote}
\usepackage{xcolor}
% ulem, not soul. `soul`'s \st and \ul are built on a character-by-character
% scan that XeLaTeX's font machinery breaks: the document fails with an
% undefined control sequence rather than degrading. ulem does the same two jobs
% and is Unicode-engine safe. `normalem` stops it hijacking \emph, which the
% transcriptions use for Grothendieck's own emphasis.
\usepackage[normalem]{ulem}
\usepackage{hyperref}
\hypersetup{colorlinks=true,linkcolor=black,urlcolor=black,pdfborder={0 0 0}}

% --- Métadonnées du lot -----------------------------------------------------
% Reprises telles quelles par le script de rendu, qui les affiche en tête de
% la vue de lecture. Elles fixent ce que le fichier prétend couvrir : sans
% elles, une transcription flotte sans point d'attache dans un fonds de seize
% mille feuillets.

\newcommand{\folder}[1]{\gdef\grothFolder{#1}}
\newcommand{\batch}[1]{\gdef\grothBatch{#1}}
\newcommand{\leaves}[2]{\gdef\grothFirst{#1}\gdef\grothLast{#2}}
\newcommand{\foldertitle}[1]{\gdef\grothFoldertitle{#1}}
\gdef\grothFolder{?}\gdef\grothBatch{?}\gdef\grothFirst{?}\gdef\grothLast{?}\gdef\grothFoldertitle{}

% --- Le résumé --------------------------------------------------------------
%
% Il ouvre la lecture modernisée, sous le titre « Résumé », et il est composé
% en petit. Ce n'est pas un ornement : c'est le seul passage du document qui
% ne soit pas des mathématiques, et un lecteur doit pouvoir le reconnaître
% avant de l'avoir lu — puis le sauter s'il connaît déjà le sujet, ou s'y
% arrêter s'il ne le connaît pas. Le filet qui le suit marque la reprise du
% corps.

\newenvironment{resume}
  {\small\setlength{\parskip}{0.45em}}
  {\par\medskip\hrule\medskip}

% --- Mots-clés ---------------------------------------------------------------
%
% En anglais, en clôture du résumé de la lecture modernisée : le vocabulaire
% moderne sous lequel chercher ce que les feuillets construisent. Le manifeste
% du site (`npm run manifest`) les extrait pour étiqueter la cote entière —
% cette ligne est la source unique des tags, il n'y a pas de fichier à côté.
\newcommand{\keywords}[1]{\par\smallskip\noindent
  {\footnotesize\textsc{Keywords} --- \textit{#1}}\par}

% --- L'appareil critique ----------------------------------------------------

% Le numéro de feuillet, en marge. C'est l'ancre qui permet de régler un
% désaccord sur une formule en regardant *un* feuillet plutôt qu'une chemise
% de six cents. Le site s'en sert aussi pour tourner les pages du fac-similé
% quand on fait défiler la transcription.
\newcommand{\leaf}[1]{%
  \marginnote{\footnotesize\color{gray!70}\textsf{#1}}%
  \label{leaf:#1}\ignorespaces}

% Illisible : on ne devine pas. Un mot inventé qui se lit comme les autres est
% le pire résultat possible.
\newcommand{\ill}{\textcolor{red!60!black}{[\,\ldots\,]}}

% Lecture douteuse : le mot est donné, et signalé comme incertain.
\newcommand{\uncertain}[1]{\uline{#1}}

% Ajout de l'éditeur — une lettre manquante, un mot sous-entendu.
\newcommand{\add}[1]{\textcolor{blue!50!black}{[#1]}}

% Biffé par Grothendieck lui-même. Ce qu'il a rayé fait partie du document :
% une rature dit souvent où la pensée a changé de direction.
\newcommand{\struck}[1]{\sout{#1}}

% Note du transcripteur, sur l'état matériel du feuillet.
\newcommand{\note}[1]{\par\smallskip{\small\color{gray!60!black}\textit{[#1]}}\par\smallskip}

% Note marginale de Grothendieck, distincte du corps du texte.
\newcommand{\marginal}[1]{\par\smallskip\noindent\hspace{1em}%
  {\small\color{gray!60!black}#1}\par\smallskip}

% --- Titre ------------------------------------------------------------------

\AtBeginDocument{%
  \begin{center}
    {\large\bfseries Fonds Alexandre Grothendieck --- cote n\textsuperscript{o}~\grothFolder}\\[2pt]
    {\itshape\grothFoldertitle}\\[4pt]
    {\small feuillets \grothFirst--\grothLast\ (lot \grothBatch)}\\[2pt]
    {\footnotesize\color{gray!60!black}Transcription établie d'après le fac-similé numérisé
     par l'Université de Montpellier.\\
     Les lectures incertaines sont soulignées, les illisibles notées [\,\ldots\,],
     les ajouts entre crochets.}
  \end{center}
  \vspace{1em}}

\endinput


\folder{151}
\batch{4}
\pages{61}{74}
\dating{[à partir de 1981-à partir de 1990]}
\watermark{Édition de démonstration\\interprétation personnelle de l'œuvre}
\foldertitle{Topos stratifié (1981) : notes manuscrites (s.d.) — lecture modernisée}

\begin{document}

\section*{Résumé}

\begin{resume}
Tout le vocabulaire est en place. Un espace $X$ est découpé en strates
$X_{i}^{*}$ indexées par un ensemble ordonné $I$ ; autour de chaque strate, un
tube $\mathcal{V}_{i,j}$ qui la voit depuis la strate voisine, et un tube épointé
$\mathcal{V}^{*}_{i,j}$ dont on a ôté la strate elle-même. Ces quatorze pages
posent la question qui reste et y répondent à moitié : \emph{quel diagramme ces
pièces forment-elles, et ce diagramme redonne-t-il l'espace ?}

Le diagramme se construit en deux temps. D'abord un recensement : quelles
inclusions existent réellement entre les pièces ? La réponse, sous des
hypothèses que l'auteur appelle raisonnables, est nette et courte — il y en a
exactement trois par couple de strates voisines, et pas une de plus. Ensuite la
combinatoire : à l'ensemble ordonné $I$ on associe un diagramme $\widetilde{I}$
dont les sommets sont les strates et les deux sortes de tubes, et dont les
flèches sont ces trois-là. Sa taille se compte : $\mathrm{card}\,I + 2e$ sommets
et $3e$ flèches, où $e$ est le nombre de couples voisins. Trois exemples sont
dessinés — une chaîne, qui donne un escalier ; un indice avec deux successeurs ;
deux indices avec un successeur commun.

Le théorème vient alors, page 68, sous le nom de \emph{théorème de recollement
des tubes} : $X$ se retrouve comme recollement du diagramme $\widetilde{I}$. Il
est donné en première version et sans démonstration — l'auteur annonce qu'elle
sera heuristique, et veut d'abord un énoncé sur les topos. On peut vérifier
qu'il est cohérent avec le début du dossier : pour $I$ à deux éléments, le
diagramme se réduit à quatre espaces et le recollement est exactement la somme
amalgamée de la section 1.

La dernière section ouvre l'appareil qui doit porter la démonstration, et c'est
de la géométrie de topos. Une stratification se transporte par image inverse le
long d'un morphisme $X' \to X$, les tubes compris ; il en résulte que le
théorème de recollement est une question \emph{locale} sur $X$, donc qu'on peut
supposer $I$ fini. Puis viennent les \emph{cribles} : les parties de $I$ closes
vers le bas, auxquelles répondent les fermés $X_{I'} = \bigcup_{i \in I'} X_{i}$,
et la correspondance respecte réunions et intersections. C'est la seconde moitié
d'un énoncé dont la section 2 avait donné la première : une stratification est
exactement une application continue de $X$ vers $I$ muni de la topologie où les
ouverts sont les parties croissantes.

Le dossier s'arrête là, au milieu d'une phrase et d'une construction. La partie
C de la section 4, qui devait généraliser les tubes mixtes aux cribles, reçoit
deux formules et rien d'autre : la page 74 est la dernière. Cette lecture
n'invente pas la suite. Ce qui se dit sans risque est que la construction
annoncée porte trois cribles là où les tubes mixtes portaient trois indices, et
que le programme énoncé à la fin de la section 3 la plaçait justement en
dernier.

\keywords{gluing theorem, sieve, base change, double mapping cylinder, codimension filtration}
\end{resume}

\pagerange{61}{61}
\section*{Le diagramme fondamental}

On rassemble les tubes en une seule pièce par degré,
\[
  \mathcal{V}_{\Delta_{1}} = \coprod_{i \leq j} \mathcal{V}_{i,j},
  \qquad
  \mathcal{V}^{*}_{\Delta_{1}} = \coprod_{i \leq j} \mathcal{V}^{*}_{i,j}
  \ \subset \ \mathcal{V}_{\Delta_{1}} \ \text{(ouvert)},
\]
ce qui donne un diagramme global reflétant, pour chaque couple $i \leq j$, le
diagramme fondamental :

\begin{tikzcd}[column sep=large]
\mathcal{V}^{*}_{i,j} \arrow[r] \arrow[d] & \mathcal{V}_{i,j} \arrow[d, dashed] & X_{i}^{*} \arrow[l] \arrow[dl, dashed] \\
X_{j}^{*} & X &
\end{tikzcd}

Les flèches en pointillé sont celles que le recollement doit fournir. Et
l'auteur note aussitôt que le regroupement global n'est probablement pas utile,
car le carré qu'il forme n'a rien de cocartésien : le diagramme en trait plein
définit bien une limite inductive, mais elle n'a guère de chances d'être
homéomorphe à $X$ — dès que $X$ est connexe, par exemple, la somme disjointe sur
\emph{tous} les couples $i \leq j$ en compte trop.

L'observation est le ressort de tout ce qui suit : ce n'est pas un carré qu'il
faut, mais un diagramme dont la forme est celle de l'ensemble ordonné $I$ tout
entier, et il faut d'abord savoir de quelles flèches il dispose.

\pagerange{62}{63}
\section*{Quelles inclusions subsistent entre les pierres}

Les $\mathcal{V}^{\,j}_{i,k}$ sont des germes de sous-espaces de $X$, entre
lesquels les seuls morphismes canoniques sont des inclusions. La question est
donc de recenser lesquelles. Pour la traiter, deux hypothèses simplificatrices
sont posées :
\[
  X_{i}^{*} \neq \emptyset \quad \text{pour tout } i,
  \qquad X_{i} \subset X_{j} \implies i \leq j,
\]
la seconde faisant de $i \mapsto X_{i}$ un plongement d'ensembles
ordonnés ;\footnote{Ce sont exactement les deux hypothèses que la section 2
s'était refusée à faire, et la section 4 dira pourquoi elle avait raison : elles
ne survivent pas au passage aux images inverses. Elles ne servent ici qu'au
recensement, qui est une question locale de géométrie et non de fonctorialité.}
et une troisième, dite « raisonnable » :
$\overline{X_{i}^{*}} = X_{i}$ pour tout $i$ — la strate est dense dans son
fermé.

Le recensement donne alors ceci. Entre les strates elles-mêmes, aucune inclusion
pour des indices distincts. Entre tubes, $\mathcal{V}_{i,j} \subset
\mathcal{V}_{i',j'}$ force $i = i'$, puisqu'il faudrait $X_{i}^{*} \subset
X_{i'}^{*}$, puis $j = j'$ puisque le tube détermine la strate voisine dont il
est le voisinage. Entre tubes épointés, $\mathcal{V}^{*}_{i,j} \subset
\mathcal{V}^{*}_{i',j'}$ donne d'abord $i \leq i'$ par passage aux adhérences,
puis, comme $\mathcal{V}^{*}_{i,j} \subset X_{j}^{*}$ et
$\mathcal{V}^{*}_{i',j'} \subset X_{j'}^{*}$ sont non vides et se rencontrent,
$j = j'$ ; et enfin $i' = i$, puisque $i \leq i' < j$ et que $j$ est un
successeur de $i$.\footnote{Cette page est surchargée de reprises et la
transcription en donne le fil plutôt que la lettre ; plusieurs des mots qui
articulent l'argument y sont illisibles. La chaîne d'implications donnée ici est
celle que les deux petits diagrammes ordonnés de la page — $i$ dominé par $i'$ et
par $j$, $i'$ dominant $j'$ — permettent de suivre, et elle aboutit à la
conclusion que la page tire. Elle est donc une reconstruction, non une lecture.}

Restent les inclusions entre types différents, et il n'y en a que trois :
\[
  X_{i}^{*} \longrightarrow \mathcal{V}_{i,j}, \qquad
  \mathcal{V}^{*}_{i,j} \longrightarrow \mathcal{V}_{i,j}, \qquad
  \mathcal{V}^{*}_{i,j} \longrightarrow X_{j}^{*}.
\]
Ce sont exactement les trois flèches en trait plein du diagramme fondamental.
Rien d'autre n'est à envisager entre ces pièces élémentaires.\footnote{Les trois
inclusions portent au manuscrit des numéros que la transcription ne lit pas ;
nous les donnons sans numéro. Le fait que le recensement retombe précisément sur
le diagramme fondamental de la page 61 est ce qui autorise la construction
suivante : le diagramme $\widetilde{I}$ n'ajoute aucune flèche qui ne soit une
inclusion effective.}

\pagerange{64}{65}
\section*{Le diagramme $\widetilde{I}$}

À l'ensemble ordonné $I$ on associe un diagramme, noté $\widetilde{I}$, dont les
sommets sont des symboles
\[
  [i], \qquad [i,j], \qquad [i,j]^{*}
  \qquad (i,j \in I,\ j \text{ successeur de } i),
\]
et dont les flèches sont de trois types, toutes décrites par

\begin{tikzcd}[column sep=large, row sep=small]
{[i,j]}^{*} \arrow[r] \arrow[d] & {[i,j]} & {[i]}^{*} \arrow[l] \\
{[j]}^{*} & &
\end{tikzcd}

Notons $e$ le nombre de \emph{drapeaux élémentaires} — les couples $i < j$ avec
$j$ successeur de $i$. Alors $\widetilde{I}$ a $\mathrm{card}\,I + 2e$ sommets et
$3e$ flèches.\footnote{Le manuscrit note $\mathcal{V}$ l'ensemble des drapeaux
élémentaires, ce qui entre en collision avec les tubes $\mathcal{V}_{i,j}$ dans
la ligne même où les deux figurent. Nous écrivons $e$ pour le cardinal.} Si $I$
a deux éléments $i < j$, on retrouve le diagramme ci-dessus lui-même.

Trois exemples sont explicités, en remplaçant les symboles par les espaces
qu'ils désignent. Pour $I = \Delta_{n} = (0 < 1 < \cdots < n)$, une chaîne, on
trouve un escalier de $n$ marches :

\begin{tikzcd}[column sep=small, row sep=small, nodes={font=\scriptsize}]
X_{0}^{*} \arrow[r] & \mathcal{V}_{0,1} & \mathcal{V}^{*}_{0,1} \arrow[l] \arrow[d] & & & & \\
& & X_{1}^{*} \arrow[r] & \mathcal{V}_{1,2} & \mathcal{V}^{*}_{1,2} \arrow[l] \arrow[d] & & \\
& & & & X_{2}^{*} \arrow[r] & \mathcal{V}_{2,3} & \mathcal{V}^{*}_{2,3} \arrow[l] \arrow[d] \\
& & & & & & \cdots \arrow[d] \\
& & & & & & X_{n}^{*}
\end{tikzcd}

Pour $I$ formé d'un indice $0$ dominé par deux indices incomparables $a$ et $b$ :

\begin{tikzcd}[column sep=small, nodes={font=\scriptsize}]
\mathcal{V}^{*}_{0,b} \arrow[r] \arrow[d] & \mathcal{V}_{0,b} & X_{0}^{*} \arrow[l] \arrow[r] & \mathcal{V}_{0,a} & \mathcal{V}^{*}_{0,a} \arrow[l] \arrow[d] \\
X_{b}^{*} & & & & X_{a}^{*}
\end{tikzcd}

Et pour $I$ formé de deux indices $a$, $b$ incomparables dominés par un même $c$ :

\begin{tikzcd}[column sep=small, nodes={font=\scriptsize}]
X_{a}^{*} \arrow[d] & & X_{b}^{*} \arrow[d] \\
\mathcal{V}_{a,c} & & \mathcal{V}_{b,c} \\
\mathcal{V}^{*}_{a,c} \arrow[u] \arrow[r] & X_{c}^{*} & \mathcal{V}^{*}_{b,c} \arrow[u] \arrow[l]
\end{tikzcd}

Dans les trois cas les comptes se vérifient : sept sommets et six flèches pour
les deux derniers, $3n+1$ sommets et $3n$ flèches pour la chaîne.

La forme de ces diagrammes mérite d'être nommée. Chaque drapeau élémentaire
$i < j$ y contribue un motif $X_{i}^{*} \to \mathcal{V}_{i,j} \leftarrow
\mathcal{V}^{*}_{i,j} \to X_{j}^{*}$, c'est-à-dire un \emph{double cylindre
d'application} entre les deux strates, passant par le tube. C'est le motif
standard par lequel un espace stratifié se reconstruit à partir de ses strates et
de ses liens, et $\widetilde{I}$ n'est que l'assemblage de ces motifs le long de
l'ensemble ordonné.\footnote{En langue d'aujourd'hui : $\widetilde{I}$ est un
modèle $1$-catégorique de la catégorie des chemins sortants, et sa colimite
homotopique est ce qui calcule le type d'homotopie de l'espace stratifié. La
formulation générale, pour les stratifications coniques, est postérieure de trente
ans ; ce qui est ici est le diagramme, et l'assertion qu'il suffit.}

\pagerange{66}{67}
\section*{La fonction dimension}

Le diagramme se simplifie si l'on dispose sur $I$ d'une fonction à valeurs
entières
\[
  d \colon I \longrightarrow \mathbb{Z},
  \qquad \text{i.e.} \quad I = \coprod_{m} I_{m},
\]
jouant le rôle d'une dimension, et satisfaisant : pour $i \leq j$,
\[
  d(j) = d(i) + 1 \iff j \text{ est un successeur de } i.
\]
On peut alors regrouper par degré :
\[
  X_{m}^{*} = \coprod_{i \in I_{m}} X_{i}^{*},
  \qquad
  \mathcal{V}_{m,m+1} = \coprod_{\substack{i \leq j \\ d(i)=m,\ d(j)=m+1}}
    \mathcal{V}_{i,j},
  \qquad
  \mathcal{V}^{*}_{m,m+1} = \coprod_{\substack{i \leq j \\ d(i)=m,\ d(j)=m+1}}
    \mathcal{V}^{*}_{i,j},
\]
d'où un diagramme $(\mathrm{Diag})_{m}$ par degré,

\begin{tikzcd}[column sep=large, row sep=small]
X_{m}^{*} \arrow[r] & \mathcal{V}_{m,m+1} & \mathcal{V}^{*}_{m,m+1} \arrow[l] \arrow[d] \\
& & X^{*}_{m+1}
\end{tikzcd}

et l'escalier obtenu en les enchaînant, $I$ étant supposé fini et les sommets
occupés par des espaces vides n'étant pas écrits.

La variante par \emph{codimension} est la même à l'orientation près : on
considère alors les $\mathcal{V}_{d,d-1}$ et $\mathcal{V}^{*}_{d,d-1}$ au lieu
des $\mathcal{V}_{d,d+1}$ et $\mathcal{V}^{*}_{d,d+1}$, et l'escalier descend au
lieu de monter, de $X_{n}^{*}$ à $X_{0}^{*}$.\footnote{Le manuscrit caractérise
la fonction codimension par « $d(j) = j - 1$ si $j$ succ. de $i$ », ce qui
confond un indice et sa dimension. La condition voulue, parallèle à celle de la
fonction dimension, est $d(j) = d(i) - 1$ ; la transcription signale le lapsus
et ne le répare pas.}

\pagerange{68}{68}
\section*{Le théorème de recollement des tubes}

\textbf{Théorème} (première version). \emph{Considérons le diagramme de type
$\widetilde{I}$ formé par les $X_{i}^{*}$, les $\mathcal{V}^{*}_{i,j}$ et les
$\mathcal{V}_{i,j}$, pour $i,j \in I$ avec $j$ successeur de $i$, et par les
trois types de flèches en trait plein du diagramme fondamental. Alors $X$
s'identifie à la limite inductive de ce diagramme.}

Le mot que le manuscrit emploie pour l'opération est illisible ; « au topos »
est écrit puis barré juste avant, ce qui indique une hésitation sur la catégorie
où prendre la limite plutôt que sur la limite
elle-même.\footnote{Nous lisons « limite inductive » parce que c'est ce que toute
la section prépare, et parce que c'est ce que le cas de la section 1 donne. La
note marginale, également reprise plusieurs fois et en partie illisible, dit
qu'il faudrait préciser en quel sens $X$ est déterminé par $\mathrm{Top}(X)$ —
c'est-à-dire prendre la limite dans les topos plutôt que dans les espaces, ce
que le mot barré suggère aussi.}

L'énoncé se vérifie sur le cas dont tout le dossier est parti. Pour $I$ à deux
éléments $0 < 1$, avec $Y = X_{0}$ et $X = X_{1}$, on a $X_{0}^{*} = Y$,
$X_{1}^{*} = X \smallsetminus Y$, $\mathcal{V}_{0,1} = \mathcal{V}_{Y,X}$ et
$\mathcal{V}^{*}_{0,1} = \mathcal{V}^{*}_{Y,X}$ ; le diagramme $\widetilde{I}$
est alors
\[
  Y \longrightarrow \mathcal{V}_{Y,X} \longleftarrow \mathcal{V}^{*}_{Y,X}
  \longrightarrow X \smallsetminus Y,
\]
et sa limite inductive est exactement la somme amalgamée de la section 1, la
flèche $Y \to \mathcal{V}_{Y,X}$ n'y ajoutant rien puisqu'elle est une
équivalence d'homotopie. Le théorème général est donc la mise en famille de ce
cas-là, sur l'ensemble ordonné $I$ tout entier.

Aucune démonstration n'est donnée. L'auteur annonce qu'elle sera heuristique, et
qu'il veut d'abord un énoncé exprimant les « topos canoniques » associés à la
situation en termes d'un diagramme de type $\widetilde{I}$ formé de topos
élémentaires. C'est l'objet de la section suivante, et le dossier s'interrompt
avant d'y revenir.\footnote{Un paragraphe barré de deux longs traits diagonaux
précède cette annonce ; il proposait de démontrer le théorème en remplaçant les
voisinages tubulaires par de vrais voisinages ouverts, faute de disposer d'une
définition assez souple des premiers. La transcription en donne le squelette. Le
détour est significatif : c'est le même que la section 1 avait pris pour rendre
la somme amalgamée élémentaire.}

\pagerange{69}{71}
\section*{4. Topos canoniques : images inverses}

Les axiomes de la section 2 sont d'abord rappelés : les $X_{i}$ fermés et la
famille localement finie ; $i \leq j \Rightarrow X_{i} \subset X_{j}$ ;
$X = \bigcup X_{i}$ ; et $X_{i} \cap X_{j} = \bigcup_{k \leq i,j} X_{k}$.

\subsection*{Transport par image inverse}

Soit $X' \to X$ un espace au-dessus de $X$. La famille des
\[
  X'_{i} = X_{i} \times_{X} X'
\]
satisfait les mêmes conditions, et le système des drapeaux se transporte :
\[
  X'_{\Delta_{r}} \simeq X_{\Delta_{r}} \times_{X} X', \qquad
  X'^{\,*}_{\Delta_{r}} \simeq X^{*}_{\Delta_{r}} \times_{X} X', \qquad
  \dot{X}'_{\Delta_{r}} = \dot{X}_{\Delta_{r}} \times_{X} X'.
\]
Ces trois identités reposent sur le fait que le changement de base commute aux
réunions de sous-objets, donc à la définition de $\dot{X}_{i}$.

C'est ici que se paie une décision de la section 2. Après image inverse, un
$X'_{i}$ peut être vide, et deux indices distincts peuvent donner le même fermé :
c'est précisément pour pouvoir prendre ces images inverses qu'il n'était pas
commode de supposer les $X_{i}$ ou les $X_{i}^{*}$ non vides, ni que
$i \mapsto X_{i}$ soit un plongement d'ensembles ordonnés dans
$\mathcal{P}(X)$.\footnote{Le manuscrit le dit explicitement, page 70, et c'est
la justification rétrospective de la remarque b) de la section 2. On notera que
les hypothèses « raisonnables » du recensement des pages 62--63 sont exactement
celles-là : elles valent pour la stratification de départ, pas pour ses images
inverses.}

Pour les tubes, le transport demande davantage. Si $X' \to X$ est une immersion
ouverte, tout va de soi. Si c'est une immersion locale, on a encore
\[
  R'_{i,j} = R_{i,j} \times_{X} X', \qquad S'_{i,j} = S_{i,j} \times_{X} X',
\]
et, dans le cas propre, des isomorphismes
\[
  \mathcal{V}'_{i,j} \simeq \mathcal{V}_{i,j} \times_{X} X',
  \qquad
  \mathcal{V}'^{\,*}_{i,j} \simeq \mathcal{V}^{*}_{i,j} \times_{X} X',
  \qquad
  \mathcal{V}'^{\,j}_{i,k} \simeq \mathcal{V}^{\,j}_{i,k} \times_{X} X'.
\]
Le manuscrit est ici plus prudent qu'ailleurs, et il a raison de l'être : la
formation d'un voisinage tubulaire ne commute pas en général à la restriction à
une partie, et il faut que $X' \to X$ soit assez proche d'une immersion ouverte
pour que le tube d'une trace soit la trace du tube.\footnote{Le haut de la page
70 est traversé par une longue note marginale diagonale plusieurs fois reprise,
qui recouvre le texte ; la transcription en donne le fil. La restriction au cas
propre, portée en marge de la formule, est de l'auteur.}

\subsection*{Une conséquence : le théorème est local}

De ces changements de base il résulte que la démonstration du théorème de
recollement est une question \emph{locale} sur $X$. On peut donc se ramener à un
ouvert au-dessus duquel la famille est finie — la locale finitude est faite pour
cela — et supposer $I$ fini. C'est le seul progrès que le dossier fasse en
direction de la démonstration, et c'en est un vrai : il ramène un énoncé sur un
ensemble ordonné quelconque à un énoncé sur un ensemble ordonné fini.

\pagerange{71}{73}
\section*{Les cribles}

Une partie $I' \subset I$ est un \emph{crible} si elle est close vers le bas :
\[
  i \leq j \in I' \implies i \in I'.
\]
On lui associe la partie fermée
\[
  X_{I'} = \bigcup_{i \in I'} X_{i}.
\]
La correspondance est monotone, et elle respecte les deux opérations :
\[
  X_{\bigcup_{\alpha} I'_{\alpha}} = \bigcup_{\alpha} X_{I'_{\alpha}},
  \qquad X_{\emptyset} = \emptyset,
  \qquad
  X_{I' \cap I''} = X_{I'} \cap X_{I''}, \qquad X_{I} = X,
\]
la troisième venant de la condition d) des axiomes — et de nulle part
ailleurs.\footnote{C'est la moitié qui manquait à la section 2. Les cribles de
$I$ sont les fermés de la topologie d'Alexandrov sur $I$, celle dont les ouverts
sont les parties croissantes ; dire que $I' \mapsto X_{I'}$ commute aux réunions
quelconques et aux intersections finies, c'est dire que $x \mapsto i(x)$ est une
application continue $X \to I$ pour cette topologie. La condition d), qui
paraissait technique, est exactement ce qui rend la stratification continue.}

Prenant $X' = X_{I'}$, on obtient sur $X'$ une stratification de type $I'$, et
les parties-cribles de $X'$ pour cette stratification sont les traces sur $X'$
des parties-cribles de $X$. Les espaces élémentaires y sont les $X_{i}^{*}$, les
$\mathcal{V}^{*}_{i,j}$ et les $\mathcal{V}_{i,j}$ pour $i,j \in I'$.

Un cas particulier vaut d'être isolé, parce qu'il justifie une désinvolture
antérieure. Soit $I_{0}$ l'ensemble des $i$ tels que $X_{i} = \emptyset$ : c'est
un crible, et tous les espaces qui lui sont associés sont vides. Le diagramme
$\widetilde{I}$, une fois les sommets vides négligés, ne change donc pas quand on
remplace $I$ par $I \smallsetminus I'_{0}$ pour un crible $I'_{0} \subset I_{0}$.
Plus généralement, pour tout crible $I'$ de $I$, le diagramme
$\widetilde{I \smallsetminus I'}$ est contenu dans $\widetilde{I}$ : le
complémentaire d'un crible est clos vers le haut, donc un couple qui y est
élémentaire l'était déjà dans $I$.

\subsection*{Ouverts entre deux cribles}

Soient maintenant deux cribles emboîtés $I'' \subset I' \subset I$, d'où
$X_{I''} \subset X_{I'}$ et la partie localement fermée
\[
  \mathcal{U}_{I',I''} = X_{I'} \smallsetminus X_{I''}.
\]
On peut y regarder la stratification de type $I' \smallsetminus I''$, définie par
\[
  X_{i',\,\mathcal{U}_{I',I''}}
  = X_{i'} \smallsetminus \bigl(X_{i'} \cap X_{I''}\bigr)
  \ \subset \ \mathcal{U}_{I',I''},
\]
dont les espaces élémentaires sont les $X_{i'}^{*}$ pour
$i' \in I' \smallsetminus I''$ et les tubes associés aux couples $(i',j')$ avec
$j'$ successeur de $i'$. Ceux-ci se comparent aux tubes de $X$ par un carré
d'inclusions,

\begin{tikzcd}[column sep=large, row sep=small]
\mathcal{V}_{i',j',X'} \arrow[r, hook] & \mathcal{V}_{i',j'} \\
\mathcal{V}^{*}_{i',j',X'} \arrow[r, hook] \arrow[u, hook] & \mathcal{V}^{*}_{i',j'} \arrow[u, hook]
\end{tikzcd}

mais — et l'avertissement est de l'auteur, il vaut d'être conservé tel quel —
il n'est pas clair en général que ces inclusions soient des équivalences
d'homotopie. Mieux vaut donc ne pas identifier trop vite les constructions faites
sur un $X_{I'}$ et celles faites sur un $\mathcal{U}_{I',I''}$.

\pagerange{74}{74}
\section*{Où le dossier s'arrête}

La partie C de la section 4 s'ouvre sur la promesse d'une famille
$\mathcal{V}^{\,I'''}_{I',I''}$, indexée par trois cribles
\[
  I',\ I''' \subset I'' \subset I,
  \qquad\text{d'où}\qquad
  X_{I'},\ X_{I'''} \subset X_{I''} \hookrightarrow X_{I} = X,
\]
et le dossier s'arrête ici. La page 74 est la dernière du fonds pour cette cote,
et rien n'indique que la suite existe ailleurs.

La construction est donc inachevée, et cette édition ne lui invente pas de fin.
Ce qui se dit sans risque tient à la notation seule : les trois cribles occupent
les places qu'occupaient les trois indices $i \leq j \leq k$ dans les tubes
mixtes $\mathcal{V}^{\,j}_{i,k}$ de la section 3, l'exposant portant l'indice
médian ; et le programme énoncé à la fin de cette même section plaçait les tubes
mixtes en dernier, « eu égard aux cas des diviseurs à croisements normaux ».
C'est cette dernière étape que la page 74 commence. Ce qu'elle aurait donné, on
ne le lit pas ici.

\end{document}
